\documentclass[11pt,a4paper]{article}

\usepackage[a4paper, margin=2cm]{geometry}
\usepackage{amsmath, amssymb}
\usepackage{booktabs}
\usepackage{xcolor}
\usepackage{longtable}
\usepackage{array}
\usepackage{parskip}
\usepackage{titlesec}
\usepackage{mdframed}

\definecolor{incgreen}{RGB}{0,140,0}
\definecolor{exred}{RGB}{180,0,0}
\definecolor{borderorange}{RGB}{200,120,0}

\newcommand{\inc}{\textbf{\textcolor{incgreen}{INCLUDE}}}
\newcommand{\exc}{\textbf{\textcolor{exred}{EXCLUDE}}}
\newcommand{\bord}{\textbf{\textcolor{borderorange}{BORDERLINE}}}

\titleformat{\section}{\large\bfseries}{}{0em}{}[\titlerule]
\titleformat{\subsection}{\normalsize\bfseries}{}{0em}{}

\pagestyle{plain}

\begin{document}

\begin{center}
  {\LARGE \textbf{Mathematics for Physics 2 --- Week 2 Review}} \\[6pt]
  {\large Sources: \textit{Memphys\_2\_week\_2.pdf}, \textit{ClassEx2 2026.pdf}} \\[4pt]
  {\normalsize Pending approval before writing to cheat sheet}
\end{center}

\vspace{1em}

\begin{mdframed}[backgroundcolor=gray!10]
\textbf{Legend:} \inc\ --- will be added to cheat sheet \quad
\exc\ --- excluded \quad \bord\ --- needs your decision
\end{mdframed}

\vspace{1em}

% ------------------------------------------------------------------
\section{Triple Products}

\begin{longtable}{p{0.5cm} p{7.5cm} p{2.2cm} p{3.5cm}}
\toprule
\# & Formula / Definition & Decision & Reason \\
\midrule
\endhead

1 & \textbf{Triple scalar product:}
\[
  V\text{ (volume)} = (\vec{a}\times\vec{b})\cdot\vec{c} = \vec{c}\cdot(\vec{a}\times\vec{b})
\]
& \inc & Named formula; geometric meaning (parallelepiped volume) \\[14pt]

2--3 & \textbf{Computation forms:}
\[
  \vec{a}\cdot(\vec{b}\times\vec{c}) =
  \begin{vmatrix} a_x & a_y & a_z \\ b_x & b_y & b_z \\ c_x & c_y & c_z \end{vmatrix}
  = \varepsilon_{ijk}\,a_i b_j c_k
\]
& \inc & Determinant form for computation; index form bridges to Levi-Civita \\[24pt]

4--5 & \textbf{Cyclic / anti-cyclic:}
\[
  \vec{a}\cdot(\vec{b}\times\vec{c}) = \vec{c}\cdot(\vec{a}\times\vec{b}) = -\vec{b}\cdot(\vec{a}\times\vec{c})
\]
& \inc & Two cyclic iterations + one anti-cyclic sign flip in a single row \\[14pt]

6 & \textbf{BAC-CAB (triple vector product):}
\[
  \vec{a}\times(\vec{b}\times\vec{c}) = \vec{b}(\vec{a}\cdot\vec{c}) - \vec{c}(\vec{a}\cdot\vec{b})
\]
& \inc & Already in Week 1 sheet; confirmed here \\[14pt]

7 & Proof of BAC-CAB via $\varepsilon_{ijk}\varepsilon_{mnk} = \delta_{im}\delta_{jn}-\delta_{in}\delta_{jm}$
& \exc & Derivation only; identity already in sheet \\[6pt]

\bottomrule
\end{longtable}

% ------------------------------------------------------------------
\section{Nabla Operator \& Gradient}

\begin{longtable}{p{0.5cm} p{7.5cm} p{2.2cm} p{3.5cm}}
\toprule
\# & Formula / Definition & Decision & Reason \\
\midrule
\endhead

8 & \textbf{Nabla operator:}
\[
  \vec{\nabla} \equiv (\partial_x,\,\partial_y,\,\partial_z) \equiv \hat{e}_i\,\partial_i
\]
& \inc & Definition of a central operator; foundation for all following entries \\[14pt]

9 & \textbf{Gradient:}
\[
  \operatorname{grad}f = \vec{\nabla}f = (\partial_x f,\,\partial_y f,\,\partial_z f) = \hat{e}_i\,\partial_i f
\]
& \inc & Core definition; scalar field in, vector field out \\[14pt]

10 & \textbf{Gradient \& differential:}
\[
  df \equiv \vec{\nabla}f\cdot d\vec{r} \approx f(\vec{r}+d\vec{r})-f(\vec{r})
\]
& \inc & Ties gradient to infinitesimal change; used constantly \\[14pt]

12 & \textbf{Key gradient result} ($r(x,y,z) = |\vec{r}(x,y,z)|$):
\[
  \vec{\nabla}r = \hat{r}
\]
& \inc & Appears in almost every physics application (central forces) \\[10pt]

13 & \textbf{Product rule:}
\[
  \vec{\nabla}(fg) = f\,\vec{\nabla}g + g\,\vec{\nabla}f
\]
& \inc & Named rule; essential for vector identities \\[14pt]

14--15 & \textbf{Chain rule (index notation --- general):}
\[
  \frac{df}{du} = (\partial_i f)\frac{dx_i}{du}
\]
\begin{tabular}{@{}l@{\;}l@{}}
  Scalar ($u=x_j$): & $\vec{\nabla}(f(g)) = f'(g)\,\vec{\nabla}g$ \\[6pt]
  Path ($u=t$):     & $\dfrac{d}{dt}f(\vec{r}(t)) = (\partial_i f)\,\dot{x}_i$
\end{tabular}
& \inc & Unified index form; both cases follow by choosing $u$ \\[28pt]

16 & Definition of partial derivative $\partial_x f := \lim_{\Delta x\to 0}\cdots$ & \exc & Calculus carry-over; assumed known \\[6pt]

17 & \textbf{Total differential (component form):}
\[
  df = \frac{\partial f}{\partial x}dx + \frac{\partial f}{\partial y}dy + \frac{\partial f}{\partial z}dz
\]
& \bord & Equivalent to entry 10 in components; may be useful or redundant \\[14pt]

18 & Derivation of gradient from Taylor expansion & \exc & Pedagogical derivation; not a formula \\[6pt]

\bottomrule
\end{longtable}

% ------------------------------------------------------------------
\section{Directional Derivative}

\begin{longtable}{p{0.5cm} p{7.5cm} p{2.2cm} p{3.5cm}}
\toprule
\# & Formula / Definition & Decision & Reason \\
\midrule
\endhead

19 & \textbf{Limit definition:}
\[
  D_{\hat{v}}f(\vec{r}) \equiv \lim_{h\to 0} \frac{f(\vec{r}+h\hat{v})-f(\vec{r})}{h}
\]
& \bord & Formal definition; exam may ask to state it \\[14pt]

20 & \textbf{Gradient form:}
\[
  D_{\hat{v}}f(\vec{r}) = \vec{\nabla}f\cdot\hat{v}
\]
& \inc & Key formula; rate of change in direction $\hat{v}$ \\[14pt]

21 & $D_{\hat{v}}f = |\vec{\nabla}f|\cos\theta$, maximised at $\theta=0$ & \inc & Gradient = direction of steepest ascent; core exam concept \\[6pt]

\bottomrule
\end{longtable}

% ------------------------------------------------------------------
\section{Curves \& Line Integrals}

\begin{longtable}{p{0.5cm} p{7.5cm} p{2.2cm} p{3.5cm}}
\toprule
\# & Formula / Definition & Decision & Reason \\
\midrule
\endhead

22 & \textbf{Parametric curve:}
\[
  \vec{r}(t) = (x(t),\,y(t),\,z(t))
\]
& \inc & Setup notation for all line-integral work \\[14pt]

23 & \textbf{Curve differential:}
\[
  d\vec{r} = \vec{r}'(t)\,dt
\]
& \inc & Key formula to convert line integrals to ordinary integrals \\[14pt]

24 & \textbf{Line integral / work:}
\[
  W = \int_C \vec{F}\cdot d\vec{r}
\]
& \inc & Central formula; definition of work as line integral \\[14pt]

25 & \textbf{Parametric form:}
\[
  W = \int_{t_1}^{t_2} \vec{F}(\vec{r}(t))\cdot\vec{r}'(t)\,dt
\]
& \inc & Computational form used in every calculation \\[14pt]

26 & Conservative force: $W=0$ on any closed path & \bord & Important concept; briefly stated, no standalone formula \\[6pt]

27 & Worked examples: axes integrals, helix, gravity & \exc & Not cheat-sheet material \\[6pt]

\bottomrule
\end{longtable}

% ------------------------------------------------------------------
\section{Gradient Theorem}

\begin{longtable}{p{0.5cm} p{7.5cm} p{2.2cm} p{3.5cm}}
\toprule
\# & Formula / Definition & Decision & Reason \\
\midrule
\endhead

28 & \textbf{Gradient theorem (multi-dimensional FTC):}
\[
  \int_{\vec{r}_1}^{\vec{r}_2} \vec{\nabla}f\cdot d\vec{r} = f(\vec{r}_2) - f(\vec{r}_1)
\]
& \inc & Path-independent; major named theorem \\[6pt]

\bottomrule
\end{longtable}

% ------------------------------------------------------------------
\section{Surfaces}

\begin{longtable}{p{0.5cm} p{7.5cm} p{2.2cm} p{3.5cm}}
\toprule
\# & Formula / Definition & Decision & Reason \\
\midrule
\endhead

29 & \textbf{Parametric surface:}
\[
  \vec{r}(u,v) = (x(u,v),\,y(u,v),\,z(u,v))
\]
& \inc & Standard definition; setup for surface integrals \\[14pt]

30 & \textbf{Surface area element:}
\[
  d\vec{S} = \left(\frac{\partial\vec{r}}{\partial u}\times\frac{\partial\vec{r}}{\partial v}\right)du\,dv
\]
& \inc & Essential formula for flux integrals \\[14pt]

31 & \textbf{Unit normal to implicit surface $f(\vec{r})=C$:}
\[
  \hat{n} = \frac{\vec{\nabla}f}{|\vec{\nabla}f|}
\]
& \inc & Standard formula; used whenever a surface normal is needed \\[14pt]

32 & Proof: $\vec{\nabla}f \perp$ tangent to surface $f=C$ & \exc & Derivation; result is entry 31 \\[6pt]

33 & Implicit surface notation: $f(x,y,z)=C$ & \bord & Context for entry 31, but just notation \\[6pt]

\bottomrule
\end{longtable}

% ------------------------------------------------------------------
\section{Applications}

\begin{longtable}{p{0.5cm} p{7.5cm} p{2.2cm} p{3.5cm}}
\toprule
\# & Formula / Definition & Decision & Reason \\
\midrule
\endhead

34 & \textbf{Extremum condition:}
\[
  \vec{\nabla}f = \vec{0}
\]
& \inc & Necessary condition for local extrema; analogue of $f'=0$ \\[10pt]

35 & Tutorial Example 1: line integral over $\vec{r}(t)=(t,t^3)$ & \exc & Worked example \\[6pt]

36 & Tutorial Example 2: cone parametrisation & \exc & Worked example \\[6pt]

\bottomrule
\end{longtable}

% ------------------------------------------------------------------
\vspace{1em}
\begin{mdframed}[backgroundcolor=orange!10]
\textbf{Items Needing Your Decision (\bord):}

\begin{tabular}{p{0.5cm} p{5cm} p{8cm}}
\toprule
\# & Item & Question \\
\midrule
17 & $df = \partial_x f\,dx + \partial_y f\,dy + \partial_z f\,dz$ & Include as alternative form of entry 10, or redundant? \\[8pt]
19 & Limit definition of directional derivative & Include in case exam asks ``state the definition''? \\[8pt]
26 & Conservative forces: $W=0$ on closed path & Include as a concept note? \\[8pt]
33 & Implicit surface notation $f(x,y,z)=C$ & Include as parenthetical alongside entry 31, or skip? \\[4pt]
\bottomrule
\end{tabular}
\end{mdframed}

\end{document}
